% Chapter 9: Final Thoughts and Future Directions.

%--------------------------------------------------------------------------
\section{Where We Started, Where We Are} \label{sec:ch9-arc}

This book began with a cave, a magic door, and Ali Baba. It ends with execution
unit measurements. That arc is deliberate: zero-knowledge proofs are one of the
few areas of computer science where the elegant intuition and the engineering
reality are separated by a very long distance, and a book that stops at the
intuition leaves you unable to build anything.

Along the way we established, by construction and measurement rather than
assertion:

\begin{itemize}
  \item A complete Groth16 verifier compiles to \textbf{570 bytes} of Plutus
        Core (Chapter~\ref{ch:offchain}).
  \item Verifying a proof costs \textbf{2{,}005{,}373{,}690 CPU steps} and
        \textbf{34{,}219 memory units}, 20\% of a transaction's CPU budget,
        about \textbf{0.147 ADA}, and that figure is \textbf{constant in the
        size of the circuit}.
  \item The cost floor for any elementary operation inside a Plutus script is
        roughly \textbf{$2\times10^5$ CPU steps} and \textbf{600 memory units},
        driven by interpreter overhead rather than operand width
        (Chapter~\ref{ch:tornado}).
  \item Those two facts together explain almost everything about what is and is
        not possible.
\end{itemize}

\begin{highlight}
\textbf{Cardano is very good at verifying zero-knowledge proofs and very bad at
computing anything ZK-adjacent on-chain.} Every conclusion in this chapter is a
corollary of that asymmetry.
\end{highlight}

%--------------------------------------------------------------------------
\section{The Central Asymmetry} \label{sec:ch9-asymmetry}

It is worth spelling out why the asymmetry exists, because it is not an accident
or an oversight.

Cardano's BLS12-381 support is implemented as \textbf{builtins}: single Plutus
Core operations backed by optimised native code, with a fixed price in the cost
model. The miller-loop builtin does an enormous amount of arithmetic and is
charged as one operation. That is why a pairing check is cheap.

Field arithmetic, by contrast, has no builtin. Aiken's \texttt{Scalar} is
ordinary Aiken code over ordinary integers, and every multiply-and-reduce is
several reduction steps in the CEK machine, each individually metered. That is
why $1{,}413$ field operations (one Poseidon hash) cost three hundred times
more than four pairings.

\begin{center}
\begin{tabular}{lrr}
  \hline
  \textbf{Operation} & \textbf{CPU steps} & \textbf{Implementation} \\
  \hline
  4 pairings + 5 decompressions (Groth16) & 2{,}005{,}373{,}690 & builtins \\
  1 Poseidon hash ($t=3$, estimated)      &   600{,}802{,}860 & Aiken code \\
  1 \texttt{blake2b\_256}                 &       398{,}720 & builtin \\
  1 field multiply-and-reduce             &       430{,}929 & Aiken code \\
  1 small integer addition                &       197{,}240 & Aiken code \\
  \hline
\end{tabular}
\end{center}

Read that table as a design principle rather than a benchmark. \textbf{On
Cardano, cryptography that has been given a builtin is nearly free, and
cryptography that has not been given a builtin is effectively unavailable.}
There is no middle ground, because the interpreter overhead per operation is
enormous relative to the arithmetic it performs. Optimising your Aiken does not
help; the only thing that helps is a builtin.

This also explains why the incremental protocol-parameter increases do not
change the picture. The transaction memory budget has gone from 11{,}250{,}000
to 14{,}000{,}000 to the current 16{,}500{,}000, about 47\% in total. A
depth-20 on-chain Poseidon insertion needs roughly 17 times the current budget.
Parameter governance moves in tens of percent; the gap is in orders of
magnitude.

%--------------------------------------------------------------------------
\section{The Missing Primitives} \label{sec:ch9-missing}

If builtins are the only lever, the natural question is which ones are missing.
In rough order of impact:

\subsection*{1. A ZK-Friendly Hash Builtin (Poseidon or Poseidon2)}

This is the single highest-leverage change available to Cardano's ZK ecosystem,
and Chapter~\ref{ch:tornado} quantified why. A Poseidon builtin would:

\begin{itemize}
  \item make on-chain Merkle tree maintenance cost roughly what
        \texttt{blake2b\_256} costs today, a factor of $\approx 1{,}500$
        improvement;
  \item keep circuits at $\approx 240$ constraints per tree level instead of
        $\approx 30{,}000$, which keeps proving keys in megabytes and proving in
        the browser;
  \item and therefore let a whole class of accumulator-based applications run
        with client-side proving and real privacy: private voting, anonymous
        credentials, membership sets, shielded pools.
\end{itemize}

\begin{highlight}
One builtin removes two of the three walls in Chapter~\ref{ch:tornado}. If you
have influence over Cardano's roadmap and you care about zero-knowledge, this is
the thing to argue for. Ethereum is having the same debate about Poseidon for
the same reasons.
\end{highlight}

Note the appropriate caution: Poseidon is young by cryptographic standards, and
arithmetisation-friendly hashes have a shorter cryptanalytic track record than
SHA-2 or Keccak. Committing it to the ledger's cost model is a long-term
commitment. Poseidon2 and its successors exist partly because the design space
is still moving. This is a real argument for care, not for inaction.

\subsection*{2. Native Field Arithmetic}

CIP-109 proposes \texttt{modularExponentiation}, which computes modular powers
and (via negative exponents) multiplicative inverses. The CIP's own
motivation notes that computing an inverse currently consumes between 5\% and
9\% of the mainnet CPU budget. It remains at \emph{Proposed} status.

That single builtin would help, but the broader need is a family of scalar field
operations priced as builtins. That is what would remove the
$2\times10^5$-steps-per-operation floor and, with it, the reason a PLONK verifier
is impractical.

\subsection*{3. A Second Proving System}

Today Groth16 is the only fully supported option, and it comes with a
circuit-specific trusted setup. The alternatives and their obstacles:

\begin{center}
\begin{tabular}{lp{7.4cm}}
  \hline
  \textbf{System} & \textbf{Status on Cardano} \\
  \hline
  \textbf{Groth16}
    & Works today. Constant 2.0B CPU verification, 192-byte proofs. One
      ceremony per circuit. \\
  \addlinespace
  \textbf{PLONK}
    & One universal, updatable ceremony for all circuits, a decisive
      advantage for any application family. Verifier needs polynomial
      commitment openings and substantial field arithmetic, which is exactly
      what Cardano prices badly. No production Aiken implementation. \\
  \addlinespace
  \textbf{Halo2}
    & No trusted setup at all, and recursion-friendly. Verifier is heavier
      still. Not viable today. \\
  \addlinespace
  \textbf{STARKs}
    & No trusted setup, hash-based rather than pairing-based. Verification is
      thousands of hash invocations, which Cardano \emph{does} have builtins
      for, making this less obviously hopeless than it first appears, though
      proof sizes strain the 16{,}384-byte transaction limit. \\
  \hline
\end{tabular}
\end{center}

\begin{remark}
The STARK row deserves a second look from anyone hunting for a research project.
Cardano's cost profile (cheap builtin hashes, ruinously expensive field
arithmetic) is close to the opposite of what a pairing-based verifier wants
and closer to what a FRI-based verifier wants. The binding constraint moves from
execution units to transaction size. Whether a Cardano-suitable STARK verifier
is possible is, as far as I know, an open and unexplored question.
\end{remark}

%--------------------------------------------------------------------------
\section{The Ecosystem, Honestly Assessed} \label{sec:ch9-ecosystem}

Cardano's ZK ecosystem in 2026 is small, real, and roughly five years behind
Ethereum's. Pretending otherwise does newcomers no favours; so does pretending
it is empty.

\begin{center}
\begin{tabular}{p{3.6cm}p{8.4cm}}
  \hline
  \textbf{Project} & \textbf{What it is} \\
  \hline
  \textbf{Janus Wallet}
    & The most complete open-source ZK application on Cardano. Password
      authentication via Groth16 with no seed phrase. Server-side proving,
      Poseidon, BLS12-381, Aiken validators. The reference for how a real ZK
      product is assembled. \\
  \addlinespace
  \textbf{aiken-zk} (Eryx)
    & Language extension and CLI that generates circuits, verifiers, and
      MeshJS redeemers from \texttt{offchain} annotations. Catalyst Fund 13.
      The most promising developer-experience work in the space, and its
      mocked Merkle hash is an honest signpost to the platform's real gap. \\
  \addlinespace
  \textbf{Modulo-P} (\texttt{ak-381}, Cardano-Semaphore)
    & A reusable Aiken Groth16 verifier and a port of Semaphore, the
      anonymous-signalling protocol. Semaphore is the same
      Merkle-tree-plus-nullifier shape as Chapter~\ref{ch:tornado}, which makes
      it the natural place to watch for progress on the accumulator problem. \\
  \addlinespace
  \textbf{cardano-foundation/bls}
    & A "gym" of BLS12-381 examples in Aiken, from the Foundation. Where to
      look for primitive-level patterns. \\
  \addlinespace
  \textbf{Midnight}
    & Cardano's ZK partner chain. Federated mainnet launched March 2026.
      Dual-ledger architecture, selective disclosure, separate token model. \\
  \hline
\end{tabular}
\end{center}

Two observations. First, almost every project on that list uses Groth16 over
BLS12-381 with Circom and snarkjs, which is why this book taught that stack: it
is not one option among many, it is the ecosystem. Second, the list is short
enough that a competent developer can read all of it in a week, which is
either intimidating or an opportunity, depending on temperament.

%--------------------------------------------------------------------------
\section{What Actually Works Today} \label{sec:ch9-works}

Chapter~\ref{ch:tornado} was a catalogue of failure, so let us be precise about
the complement. The test is the one from Section~\ref{sec:ch8-verdict}: does the
proof concern data the chain already holds, or state the chain must maintain?

\subsection*{Authentication and Access Control}

Prove knowledge of a credential without revealing it. The commitment sits in a
datum; the validator reads it. This is Chapter~\ref{ch:password}, and Janus
Wallet is the production version. \textbf{Works today, entirely.}

\subsection*{Off-Chain Computation Compression}

Move an expensive computation off-chain and prove it was performed correctly.
Because verification is constant at 2.0B CPU steps regardless of circuit size,
this trade is spectacular: an arbitrarily large computation, verified for 0.147
ADA. This is precisely what \texttt{aiken-zk}'s \texttt{offchain} blocks are
for. \textbf{Works today, and is underexploited.}

\begin{highlight}
If you are looking for a ZK project on Cardano with a favourable
effort-to-value ratio today, this is it. The constant verification cost means
the more expensive your computation, the better the trade looks, and no
accumulator is involved, so none of Chapter~\ref{ch:tornado}'s walls apply.
\end{highlight}

\subsection*{Selective Disclosure and Attestations}

Prove you are over 18, or resident in a jurisdiction, or accredited, without
revealing the underlying document. The issuer's public key or credential root
lives in a datum or a reference input; the chain maintains nothing. \textbf{Works
today.}

\subsection*{Private Voting with Off-Chain Roll-Up}

Prove membership in an eligible-voter set and cast a vote nobody can link to
you. This works if the membership root is fixed at the start of the vote:
published once, then read-only. It hits Chapter~\ref{ch:tornado}'s walls only if
the register must be updated on-chain during the vote. \textbf{Works today, with
that design constraint.}

\subsection*{Cross-Chain and Light-Client Proofs}

Prove that a block header chain or state transition on another chain is valid.
Verification is one constant-cost check. The interesting nuance is that the
circuit must implement the \emph{other} chain's hash function, which is usually
Keccak or SHA-256, expensive in-circuit, exactly the Wall Two trade-off, but
here there is no on-chain accumulator forcing the issue. \textbf{Works today, at
the cost of large circuits.}

\subsection*{The Pattern}

\begin{center}
\begin{tabular}{p{5.6cm}p{6.4cm}}
  \hline
  \textbf{Works today} & \textbf{Blocked today} \\
  \hline
  Statement about a fixed public value in a datum
    & Statement about an accumulator the chain must update \\
  \addlinespace
  One proof, one transaction, no shared state
    & Many users contending for one state UTxO \\
  \addlinespace
  Root published once, read many times
    & Root recomputed on-chain per transaction \\
  \addlinespace
  Circuit-expensive hashes acceptable
    & Requires a hash cheap both on-chain and in-circuit \\
  \hline
\end{tabular}
\end{center}

%--------------------------------------------------------------------------
\section{Layer 2 and the Midnight Answer} \label{sec:ch9-l2}

Cardano's institutional answer to the applications in the right-hand column is
not to reshape the eUTxO ledger. It is to move them.

\textbf{Hydra} removes contention by moving state into a head with a fixed
participant set. That is genuinely useful for throughput and useless for
anonymity (a known participant set is the opposite of an anonymity set), but
it is a strong fit for ZK applications where the parties are known and the
\emph{data} is what needs hiding.

\textbf{Midnight} is the more direct answer: a partner chain designed from the
start for ZK smart contracts, with a dual-ledger architecture separating public
from private data, and selective disclosure as a first-class feature rather than
a bolt-on. Its federated mainnet launched in March 2026.

The strategic logic is coherent, and worth understanding rather than merely
noting. Cardano's L1 optimises for deterministic, locally verifiable, cheaply
priced validation. Those properties are why the eUTxO model has no shared
mutable state, and the absence of shared mutable state is exactly Wall Three.
You cannot have both. So the L1 keeps its guarantees and hosts ZK
\emph{verification}, and a purpose-built chain hosts ZK \emph{state}.

\begin{remark}
Be alert to the trade this makes on your behalf. Midnight's selective disclosure
is a deliberate design position: privacy that can be lifted under defined
conditions. That is a genuine answer to the regulatory problem that opened
Chapter~\ref{ch:tornado}, and it is also, unavoidably, weaker than
unconditional privacy. Whether that is the right trade depends on the
application and on who you think should hold the key. It is not a technical
question, and it should not be presented as one.
\end{remark}

%--------------------------------------------------------------------------
\section{Ethics, Regulation, and Building Anyway} \label{sec:ch9-ethics}

This book contains a chapter analysing a sanctioned protocol. It seems worth
being direct about how to think about that.

Zero-knowledge proofs are unusually dual-use. The same construction that lets a
refugee prove eligibility for aid without revealing their identity lets someone
launder proceeds of crime. The mathematics does not distinguish. Neither, in the
Tornado Cash case, did the enforcement action: OFAC sanctioned the smart
contract addresses themselves, and a developer was prosecuted over software he
did not control at the time of the conduct at issue. Those actions remain
contested, and a fuller reckoning belongs to lawyers and courts rather than to
this book.

What belongs here are three practical positions.

\textbf{Privacy is the default in every other part of life, and its absence on
public blockchains is an artefact rather than a virtue.} A blockchain where
every salary payment, medical invoice, and political donation is permanently
public to everyone is not a neutral design; it is a surveillance system that
happens to be decentralised. Work on privacy needs no special justification.

\textbf{Threat modelling is part of engineering.} Chapter~\ref{ch:offchain}'s
point about server-side proving is the concrete form of this: a system marketed
as zero-knowledge whose prover runs on someone else's server provides no privacy
against that someone. If you build ZK systems, state clearly what is hidden,
from whom, and under what assumptions. Users cannot evaluate a claim like "uses
zero-knowledge proofs".

\textbf{Selective disclosure is a legitimate engineering answer, not a
capitulation.} Systems where the user can prove specific facts to specific
counterparties (rather than either revealing everything or nothing) cover
most real requirements, and are the direction Cardano and Midnight have chosen.
Absolute anonymity is a narrower need than the discourse suggests.

\begin{highlight}
The interesting question is not "should privacy exist" but "who holds the
disclosure key, and under what process". That is a governance question. Building
systems that make it an explicit, auditable governance question, rather than
an implicit property of whoever runs the server, is a contribution.
\end{highlight}

%--------------------------------------------------------------------------
\section{What to Build Next} \index{next steps} \label{sec:ch9-next}

Concrete suggestions, ordered by ratio of value to difficulty, for anyone who
finished this book and wants to contribute.

\begin{enumerate}
  \item \textbf{Add replay protection to the Chapter~\ref{ch:password}
        contract.} Add a nonce to the datum and the circuit's public inputs,
        require the continuing output to increment it. This is the single most
        important missing feature, it is a genuine exercise, and it takes an
        afternoon.

  \item \textbf{Implement Poseidon in Aiken and publish the measurements.} Yes,
        it will be too expensive to use; that is the point. Chapter
        \ref{ch:tornado}'s figures are estimates extrapolated from measured
        primitives. A real implementation with real numbers would turn them into
        data, and data is what a builtin proposal needs.

  \item \textbf{Build something using off-chain computation compression.} Take
        an expensive on-chain computation, move it into a circuit, verify the
        proof for a constant 0.147 ADA. No accumulator, no contention, no walls.
        The best available ratio of effort to result on Cardano today.

  \item \textbf{Write the missing tooling.} A library that generates Aiken
        types and MeshJS encoders from a \texttt{verification\_key.json} would
        eliminate most of the failure modes in
        Section~\ref{sec:ch7-failures}. The bugs in this book's tables are
        encoding bugs, and encoding bugs are what code generation exists to
        prevent.

  \item \textbf{Investigate a STARK verifier for Aiken.} Speculative and
        possibly a dead end. Also the one item on this list where Cardano's
        peculiar cost profile (cheap builtin hashes, expensive field
        arithmetic) might turn out to be an advantage rather than a handicap.

  \item \textbf{Argue for a Poseidon builtin, with numbers.} The measurements in
        Chapter~\ref{ch:tornado} exist to be reused. Cardano's governance
        process responds to concrete cases.
\end{enumerate}

%--------------------------------------------------------------------------
\section{Where to Continue} \label{sec:ch9-resources}

\begin{center}
\begin{tabular}{p{5.1cm}p{6.9cm}}
  \hline
  \textbf{Resource} & \textbf{Why} \\
  \hline
  \texttt{aiken-lang.org}
    & Language reference and the \texttt{stdlib} \texttt{bls12\_381} modules.
      Read the \texttt{scalar} source; it explains a great deal. \\
  \addlinespace
  \texttt{docs.circom.io}
    & The circuit language, and the underlying R1CS model. \\
  \addlinespace
  \texttt{github.com/iden3/circomlib}
    & The primitive library. Read \texttt{poseidon.circom} against
      Chapter~\ref{ch:tornado}'s round counts. \\
  \addlinespace
  \texttt{github.com/leobel/janus-wallet}
    & A complete ZK product. Read the circuit and the Aiken validators
      together. \\
  \addlinespace
  \texttt{github.com/eryxcoop/cardano-zk-aiken}
    & Where Cardano's ZK developer experience is being built. \\
  \addlinespace
  \texttt{cips.cardano.org}
    & CIP-101 (keccak256), CIP-109 (modular exponentiation), CIP-127
      (ripemd-160). How primitives actually arrive. \\
  \addlinespace
  \texttt{zkproof.org}
    & Standards work, and a reading path into the underlying literature. \\
  \hline
\end{tabular}
\end{center}

And the two papers worth reading in full rather than summarised: Goldwasser,
Micali and Rackoff's 1985 paper for what a zero-knowledge proof \emph{is}, and
Groth's 2016 paper for the construction whose verification equation we
implemented in Chapter~\ref{ch:password}.

%--------------------------------------------------------------------------
\section{Closing} \label{sec:ch9-closing}

This book set out to make zero-knowledge proofs on Cardano buildable rather than
merely admirable. Nine chapters later, the honest summary is two sentences.

\begin{outline}
Cardano can verify a zero-knowledge proof about any computation, of any size,
for about 0.147 ADA and 570 bytes of on-chain code. It cannot yet maintain the
kind of cryptographic state that the most ambitious ZK applications need, and
the reason is a missing hash builtin and a ledger model that deliberately has no
shared mutable state.
\end{outline}

The first sentence is a remarkable capability that very few developers on
Cardano are currently using. The second is a real limitation with a known,
tractable, single-item fix at the top of its list.

Neither sentence was findable in one place when this project started, which is
why the project existed. Both are now measured, reproducible from the repository,
and free to reuse.

\begin{highlight}
Everything in this book (the LaTeX source, the Circom circuits, the Aiken
validators, the benchmarks, the measurements) is open source under AFL-3.0.
Corrections are more welcome than praise; the numbers in Chapter
\ref{ch:tornado} in particular are extrapolations that deserve to be replaced by
someone's real implementation.

Go build something. Publish what breaks.
\end{highlight}
